% ======================================================
%  "Jak korzystać z tej książki" -- the one page that makes the format work.
%
%  It has to fit on ONE page: a reader who has to turn over to finish learning
%  how to use the book has already been shown the thing the method depends on
%  not seeing.
%
%  NOTE THE WORDING. Every instruction here is true under any pagination. The
%  book must never say "zanim odwrócisz stronę" -- where a page breaks is a
%  property of the format, not of the text, so an instruction that names a page
%  turn is true in one build and false in another. "Zanim przeczytasz dalej"
%  and "zakryj stronę poniżej kreski" both name the reader's hand, which is
%  what the method actually asks for, and are true everywhere.
% ======================================================

\chapter*{Jak korzystać z tej książki}
\addcontentsline{toc}{chapter}{Jak korzystać z tej książki}

Każde zadanie zajmuje jedną \emph{ramkę} — kawałek strony między dwiema
cienkimi kreskami, z numerem na zewnętrznym marginesie. Odpowiedź nie stoi pod
zadaniem: otwiera ramkę \emph{następną}. Dzięki temu możesz zakryć dłonią
stronę poniżej kreski i pracować, nie widząc rozwiązania.

\begin{enumerate}
  \item Zakryj dłonią stronę poniżej zadania.
  \item Uruchom stoper.
  \item Rozwiąż zadanie \emph{w pamięci}. Bez kalkulatora, bez kartki.
  \item Zatrzymaj stoper i porównaj swój czas z limitem podanym przy zadaniu.
  \item Dopiero teraz czytaj dalej — odpowiedź czeka w kolejnej ramce.
\end{enumerate}

\noindent
Nie chodzi o to, żeby zawsze zmieścić się w limicie. Chodzi o to, żeby
\emph{spróbować odpowiedzieć, zanim zobaczysz rozwiązanie} — także wtedy, gdy
nie jesteś pewien. Pomyłka, którą sam poprawisz, zostaje w pamięci lepiej niż
poprawna odpowiedź, którą tylko przeczytałeś.

\section*{Gwiazdki i czas}

Gwiazdki przy zadaniu mówią, ile powinno zająć:

\begin{center}
  \begin{tabular}{@{}c@{\hspace{1.5em}}l@{\hspace{1.5em}}l@{}}
    \DifficultyStars{1} & łatwe   & do 1 minuty \\[3pt]
    \DifficultyStars{2} & średnie & 1--3 minuty \\[3pt]
    \DifficultyStars{3} & trudne  & 3--5 minut \\
  \end{tabular}
\end{center}

\noindent
Obok gwiazdek stoi kategoria zadania i limit czasu:

% On its own centred line rather than run into the sentence. Inline, the two
% badges are unbreakable boxes that TeX will happily split across a line break,
% so the category ends one line and the clock starts the next -- which reads as
% two unrelated things instead of one example of a frame's header.
\begin{center}
  \CategoryBadge[ArithColor!20]{Dodawanie}\hspace{0.6em}\TimerBadge{30 s}
\end{center}

\section*{Pięć obszarów}

Rozdziały są niezależne — możesz zaczynać od dowolnego i wracać do nich w
dowolnej kolejności. Na otwarciu każdego rozdziału jest napisane, ile ma
ramek, więc wiesz, na co się piszesz, zanim zaczniesz.
